% Chapitre « Données et variables » (partie 3, jalon J4).
% Sources : docs/MODELE_MATHEMATIQUE.md, ADR-0008 à 0012, 0023, 0028 à 0033, et le code de
% src/foot_predictor/features/ (elo.py, rolling.py, xg_proxy.py, rest.py, sources.py, seal.py).

\chapter{Données et variables}
\label{chap:donnees}

Ce chapitre décrit d'où viennent les données, la règle temporelle qui s'impose à toutes les
variables, puis les variables du MVP : l'Elo (groupe G1), les moyennes glissantes et le
\og xG estimé \fg{} (groupe G2), le calendrier et le contexte (groupes G3 et G0). Il se termine
par le scellé, qui protège la saison de test. Les paramètres cités sont ceux figés le
29 septembre 2026.

% ---------------------------------------------------------------------------------------------
\section{Sources et référentiel}
\label{sec:sources}

Les données brutes sont des fichiers immuables : les réponses d'API-FOOTBALL (de 2010 à 2026
selon les championnats) et les fichiers CSV de football-data (de 2000-01 à 2026-27). Une
commande reconstruit entièrement le référentiel \code{staging} à partir du brut (ADR-0008) :
\begin{itemize}
  \item les équipes, les matchs et les joueurs sont identifiés par les identifiants
        d'API-FOOTBALL ;
  \item un match de football-data est apparié au match API de même équipe à domicile, même équipe
        à l'extérieur, à un jour près ; hors de la couverture de l'API, football-data crée le
        match (origine \og hors API \fg{}) ;
  \item les tirs et tirs cadrés de chaque équipe viennent de football-data (ADR-0029).
\end{itemize}
La \emph{cible} est le nombre de buts de chaque équipe au temps réglementaire (ADR-0009). Les matchs
annulés, abandonnés ou sur tapis vert sont exclus : ils ne donnent ni ligne ni historique.

Le périmètre du MVP compte dix championnats (les cinq premières divisions européennes et leurs
deuxièmes divisions), de 2000-01 à 2024-25 : \(98\,049\) matchs de saison régulière, donc
\(196\,098\) lignes \((m, e)\), une par match \(m\) et par équipe \(e\).

% ---------------------------------------------------------------------------------------------
\section{Règle temporelle}
\label{sec:regle}

Notons \(j_m\) le jour d'un match (date UTC du coup d'envoi ; date de football-data pour un match
hors API) et \(\hist\) l'historique des matchs terminés. Pour un jour \(J\), l'historique
disponible est
\[
  \hist_{<J} = \{\, m \in \hist : j_m < J \,\}.
\]
Toute variable \(z\) d'une ligne \((m, e)\) jouée le jour \(J = j_m\) est une fonction de cet
historique seulement : \(z_{m,e} = f(\hist_{<J}, e)\). Deux conséquences, vérifiées par des tests :
\begin{enumerate}
  \item modifier le résultat du match \(m\) ne change pas la ligne de \(m\) ; deux matchs du même
        jour ne s'influencent pas ;
  \item \emph{invariance à la date de coupe} : pour toute date \(d\) et tout \(J < d\),
        \(\big(\hist_{<d}\big)_{<J} = \hist_{<J}\), donc les lignes calculées sur l'historique
        tronqué à \(d\) sont identiques à celles calculées sur l'historique complet. Sur les
        données réelles, l'égalité est exacte, au bit près, aux dates de coupe du 1\textsuperscript{er}
        janvier 2019 et du 1\textsuperscript{er} mars 2023.
\end{enumerate}

% ---------------------------------------------------------------------------------------------
\section{Elo (groupe G1)}
\label{sec:elo}

Chaque pays forme une \emph{échelle} : sa première et sa deuxième division, liées par la montée et
la descente. Il n'y a aucune échelle commune entre pays (ADR-0031).

\paragraph{Espérance.} Pour un match entre une équipe à domicile de note \(R_d\) et une équipe à
l'extérieur de note \(R_x\), l'espérance de résultat de l'équipe à domicile est
\[
  E = \frac{1}{1 + 10^{-(R_d + H - R_x)/400}},
\]
où \(H\) est l'avantage du terrain, en points. C'est la probabilité de victoire plus la moitié de
celle du nul.

\paragraph{Mise à jour.} Avec le résultat \(S \in \{1, \tfrac12, 0\}\) et l'écart de buts
\(N = |b_d - b_x|\),
\[
  \Delta R_d = K \, G(N) \, (S - E), \qquad \Delta R_x = -\Delta R_d,
  \qquad
  G(N) =
  \begin{cases}
    1 & \text{si } N \le 1,\\
    3/2 & \text{si } N = 2,\\
    (11 + N)/8 & \text{si } N \ge 3.
  \end{cases}
\]
La somme des notes d'une échelle est donc conservée par chaque match. Les variations d'un même
jour sont calculées avec les notes du début du jour et appliquées ensemble à la fin du jour : la
note d'avant-match ne dépend que de \(\hist_{<J}\).

\paragraph{Intersaison.} Au premier match d'une nouvelle saison, une équipe présente dans l'échelle
la saison précédente revient vers \(m_\ell\), la note moyenne de fin de saison précédente des
équipes de sa division \(\ell\) de la nouvelle saison :
\[
  R \leftarrow m_\ell + (1 - r)\,(R - m_\ell).
\]
Le niveau d'une équipe dans une saison se lit sur la saison régulière (un barrage ne compte pas).
Une équipe sans match dans l'échelle la saison précédente reçoit la moyenne des trois plus basses
notes de fin de saison de sa division ; la première saison (2000-01), 1500 en première division et
\(1500 - \Delta\) en deuxième.

\paragraph{Réglage.} Les paramètres \((K, H, r, G, \Delta)\) sont choisis sur une grille fixée à
l'avance (576 combinaisons), évaluée sur les saisons 2005-06 à 2014-15 seulement, avant le début
de l'apprentissage. Pour chaque combinaison, on ajuste une logistique ordonnée sur
\(x = (R_d + H - R_x)/400\) :
\[
  P(\text{extérieur}) = \sigma(c_1 - \beta x),\quad
  P(\text{nul}) = \sigma(c_2 - \beta x) - \sigma(c_1 - \beta x),\quad
  P(\text{domicile}) = 1 - \sigma(c_2 - \beta x),
\]
avec \(\sigma(u) = 1/(1+e^{-u})\), et l'on retient la plus petite log-loss moyenne. Retenu :
\(K = 15\), \(H = 40\), \(r = 0\), multiplicateur \(G\), \(\Delta = 100\) ; log-loss
\(0{,}9936\) sur \(18\,259\) matchs, contre \(1{,}0622\) pour les seules fréquences de 1, N et 2.

% ---------------------------------------------------------------------------------------------
\section{Moyennes glissantes (groupe G2)}
\label{sec:glissants}

\paragraph{Définition.} Soit \(x_i\) une quantité mesurée dans un match de championnat \(i\) de
l'équipe (buts marqués, buts encaissés, \(\mathrm{xg}\) estimé pour ou contre), joué le jour
\(j_i\). Pour le match du jour \(J\), avec une demi-vie \(h\) (en jours) et une ancienneté maximale
\(A = 730\) jours, les poids sont
\[
  w_i(J) = 2^{-(J - j_i)/h} \quad \text{si } 1 \le J - j_i \le A, \qquad w_i(J) = 0 \text{ sinon},
\]
et, avec \(W = \sum_i w_i\) et un poids a priori \(k = 3\),
\[
  \hat{x} = \frac{\sum_i w_i\, x_i + k\, \mu}{W + k},
\]
où \(\mu\) est la moyenne par équipe et par match de la même quantité dans le championnat du
match, sur les \(A\) jours avant \(J\). Si \(W = 0\) (aucun historique), \(\hat{x}\) reste
\emph{vide} : on ne remplace jamais un manque par 0 ni par \(\mu\). La somme \(W\) est elle-même
une variable : elle mesure la quantité d'information. Trois demi-vies candidates sont calculées
(\(h \in \{60, 120, 240\}\)) ; le choix se fait en partie 4, par validation interne.

\paragraph{Propriété.} Posons \(\lambda = W/(W + k) \in [0, 1)\) et
\(\bar{x}_w = \sum_i w_i x_i / W\). Alors
\[
  \hat{x} = \lambda\, \bar{x}_w + (1 - \lambda)\, \mu,
\]
une combinaison convexe entre la moyenne pondérée de l'équipe et celle du championnat. Quand
\(W \to \infty\), \(\hat{x} \to \bar{x}_w\) ; quand \(W \to 0^{+}\), \(\hat{x} \to \mu\). Pour une
équipe qui joue une fois par semaine depuis longtemps, sans ancienneté maximale,
\(W \approx 1/(1 - 2^{-7/h})\), soit environ 12,9 matchs pour \(h = 60\) et 25,2 pour \(h = 120\) :
\(k = 3\) pèse alors peu.

\paragraph{Pourquoi en jours.} Une décroissance par nombre de matchs donnerait le même poids au
dernier match de mai et au match de la semaine précédente, quelle que soit la pause d'été. En
jours, l'été compte : un match de mai pèse moins en août. Il n'y a pas de remise à zéro à chaque
saison (rapport de cadrage, H.4).

\paragraph{Calcul exact.} Le poids se sépare : \(2^{-(J - j)/h} = 2^{-J/h} \cdot 2^{j/h}\). Si
\(C(n) = \sum_{i \le n} 2^{j_i/h} x_i\) est la somme cumulée dans l'ordre chronologique, la somme
sur la fenêtre vaut \(2^{-J/h}\,\big(C(\text{fin}) - C(\text{début})\big)\). Un historique tronqué
à une date en est exactement le préfixe : c'est ce qui rend l'invariance à la date de coupe exacte.

% ---------------------------------------------------------------------------------------------
\section{xG estimé à partir des tirs}
\label{sec:xgproxy}

L'xG d'API-FOOTBALL n'existe pas avant 2022-23 (ADR-0023). Le MVP le remplace par une combinaison
linéaire fixe des tirs cadrés \(s_c\) et non cadrés \(s_n\) (tirs contrés compris) de l'équipe :
\[
  \mathrm{xg\_proxy} = a\, s_c + b\, s_n,
  \qquad
  (a, b) = \operatorname*{arg\,min}_{a \ge 0,\ b \ge 0}\ \sum_{(m,e)} \big(y_{m,e} - a\, s_{c} - b\, s_{n}\big)^2,
\]
sur les \(40\,394\) lignes des saisons 2015-16 à 2020-21, avant tous les plis de validation. Pas de
constante (zéro tir donne zéro) et des coefficients positifs ou nuls (un tir ne retire pas de but
attendu), comme pour un vrai xG, somme de probabilités de but par tir.

Sans la contrainte, on trouve \(b = -0{,}0126\). La contrainte \(b \ge 0\) est donc active, et les
conditions d'optimalité donnent \(b = 0\) et \(a = \sum y\, s_c / \sum s_c^2\) :
\[
  \mathrm{xg\_proxy} = 0{,}3033\, s_c.
\]
Contrôle sur 2022-23 à 2024-25, contre l'xG d'API-FOOTBALL, par équipe et par match : corrélation
\(0{,}685\), écart moyen absolu \(0{,}495\). Limite : les tirs de football-data changent de
définition pour la Serie A de 2018-19 à 2020-21, ce qui y surestime l'\(\mathrm{xg\_proxy}\)
d'environ \(0{,}27\) par équipe et par match (ADR-0029, ADR-0032).

% ---------------------------------------------------------------------------------------------
\section{Calendrier et contexte (groupes G3 et G0)}
\label{sec:calendrier}

Sur les matchs terminés avant \(J\), toutes compétitions présentes dans \code{staging} :
\begin{itemize}
  \item jours de repos \(J - j_{\text{précédent}}\) ;
  \item nombre de matchs du jour \(J - 14\) au jour \(J - 1\) ;
  \item présence d'un match de coupe d'Europe du jour \(J - 4\) au jour \(J - 1\).
\end{itemize}
Ces variables sont rétrospectives : la date du prochain match n'est jamais utilisée. Un indicateur
dit si toutes les coupes du pays sont connues cette saison (depuis 2015-16, 2016-17 pour l'Italie,
2018-19 pour l'Espagne). Le contexte (G0) comprend le lieu (domicile ou extérieur), le championnat,
la saison et un indicateur de huis clos, établi par périodes sourcées ; une période incertaine vaut 0
(ADR-0033).

% ---------------------------------------------------------------------------------------------
\section{Le scellé}
\label{sec:scelle}

La période de développement est \(\mathcal{D} = \{\, m : j_m < \text{1\textsuperscript{er} juillet 2025} \,\}\)
(ADR-0012). Les matchs hors de \(\mathcal{D}\) forment la saison de test, utilisée une seule fois
par version du modèle. Techniquement (ADR-0028) :
\begin{itemize}
  \item une seule date dans le code, et une seule porte de lecture des matchs, qui filtre dans la
        requête SQL (\(j_m < \) 1\textsuperscript{er} juillet 2025), puis vérifie le résultat ;
  \item lire un match scellé exige une option explicite, et chaque usage ajoute une ligne à un
        journal versionné ;
  \item un test d'architecture interdit tout autre accès aux matchs dans les variables et les
        notebooks.
\end{itemize}
Le jeu de données du 29 septembre 2026 ne contient aucun match scellé, et le journal des tests
scellés est vide.
